\section*{Step 33: F24 Type-Uniqueness of the QM--GR Role-Split}

\paragraph{Finite data.}
Let \(H\) be the eight-state finite carrier over modes
\[
  (d_0\ {\rm density}, d_1\ {\rm phase}, d_2\ {\rm transport}, d_3\ {\rm curvature}).
\]
Let \(q_{\rm QM}=(d_0,d_1,d_2)\), let \(s=d_3\), and let
\[
  \pi_L=(d_0,d_1,d_2,d_3).
\]
The maps \(a\pi_L=q_{\rm QM}\) and \(b\pi_L=s\) are the actual maps used in Steps 17--18 and structuralized in Step 21.

\paragraph{Computation.}
The script \texttt{type\_uniqueness\_step33.py} evaluates all eight structural F24 predicates on \(L\).  The predicate table records exactly one firing predicate:
\[
  {\rm BridgeMediatedRole}(L)=\top,
\]
with all seven remaining families false.  The coarsening search checks every quotient obtained by retaining a subset of the \(q_{\rm QM}\) modes; every such coarsening leaves \(O_s\) nonempty, with minimum obstruction count \(3\).

\paragraph{Verdict.}
Within the declared finite grammar, the QM--GR role-split has a unique F24 resolution type:
\[
  \boxed{{\rm BridgeMediatedRole}.}
\]
The hidden-latent, coarsened-only, budget-only, scoped-only, memory, outside-scope, and blocked-nonclosure alternatives are computed-excluded as full reconciliation types for this \(L\).

\paragraph{Nonclaim.}
This is type-uniqueness only.  It does not prove instance-uniqueness among all possible mediated objects, and it does not certify frame transfer beyond the finite carrier.
